% Generated draft: requires mathematical and source review.
\documentclass[11pt]{article}
\usepackage[margin=25mm]{geometry}
\usepackage{fontspec}
\setmainfont{Latin Modern Roman}
\usepackage{amsmath,amssymb,mathtools,unicode-math}
\setmathfont{Latin Modern Math}
\usepackage{enumitem,needspace,xurl,hyperref}
\setlength{\parindent}{0pt}
\setlength{\parskip}{5pt}
\setlength{\emergencystretch}{3em}
\newcommand{\N}{\mathbb{N}}
\newcommand{\Z}{\mathbb{Z}}
\newcommand{\Q}{\mathbb{Q}}
\newcommand{\R}{\mathbb{R}}
\newcommand{\C}{\mathbb{C}}
\newcommand{\F}{\mathbb{F}}
\newcommand{\A}{\mathbb{A}}
\newcommand{\PP}{\mathbb{P}}
\newcommand{\E}{\mathbb{E}}
\newcommand{\eps}{\varepsilon}
\newcommand{\dd}{\,\mathrm{d}}
\begin{document}
\section*{UnsolvedMath ID: 30002985}
\textbf{Source problem code:} catalogue entry
\subsection{Definitions and mathematical statement}
The supplied record poses two questions. Its base field, topology, and meaning of ``fake K3 surface'' are not specified; the following makes the conventional interpretation of the first question explicit while retaining those limitations.

For an intended algebraically closed base field \(k\), let \(M_g\) denote the moduli space of smooth, projective, geometrically connected curves of genus \(g\) over \(k\). Under the projective interpretation, a K3 surface is a smooth, projective, geometrically connected surface \(S\) with \(\omega_S\cong\mathcal O_S\) and \(H^1(S,\mathcal O_S)=0\). Define
\[
(K3)_g=\bigl\{[C]\in M_g(k):\text{there exist a K3 surface }S/k\text{ and a closed immersion }C\hookrightarrow S\bigr\}.
\]
No primitivity, ampleness, or polarization-degree restriction is stated in the supplied problem.

\begin{enumerate}
\item For each integer \(g\geq12\), is \((K3)_g\) closed in \(M_g\)? If ``closed'' means Zariski closed, the question is whether this set is the set of \(k\)-points of a Zariski-closed subset of \(M_g\), with closure taken inside the moduli space of smooth curves. The source does not specify the field or its characteristic, nor explicitly identify the topology; this formulation is conditional on these conventions and does not assert a question uniformly over all fields.
\item Can Lazarsfeld's proof of the Gieseker--Petri theorem be generalized to curves on fake K3 surfaces? For reference, for a smooth projective curve \(C\) and a line bundle \(L\) on \(C\), the Petri multiplication map is
\[
\mu_{C,L}:H^0(C,L)\otimes_k H^0(C,\omega_C\otimes L^{-1})\longrightarrow H^0(C,\omega_C).
\]
The Petri condition is injectivity of \(\mu_{C,L}\) for every line bundle \(L\). The supplied record does not define fake K3 surfaces, identify the allowed curves or linear systems on them, or specify whether a general or every curve is intended. Consequently, this second question remains a request for a proof generalization with unresolved hypotheses; it cannot be replaced by a precise universal Petri assertion from the supplied information.
\end{enumerate}
\subsection{Short English statement}
For each \(g\geq12\), ask whether the smooth genus-\(g\) curves that embed in K3 surfaces form a closed locus in \(M_g\). Separately, ask whether Lazarsfeld's Gieseker--Petri argument extends to curves on fake K3 surfaces. The field, topology, meaning of fake K3, and hypotheses of the proposed extension require clarification.
\subsection{Sources}
\begin{itemize}
\item Local UnsolvedMath registry, supplied complete record: permanent ID \(30002985\), source problem code \texttt{catalogue entry}, title ``Closedness of the K3 Locus and Fake K3 Surfaces,'' collection ``Oberwolfach Reports Open Problems.'' No individual proposer is supplied; the proposed year is recorded as 2015.
\item Supplied report citation: \emph{Mini-Workshop: Singular Curves on \(K3\) Surfaces and Hyperkähler Manifolds} (2016), \emph{Oberwolfach Reports}, pp.\ 2939--2967, DOI \texttt{10.4171/owr/2015/51}. \url{https://doi.org/10.4171/owr/2015/51}. The supplied report title uses 2016, whereas the proposed year, year tag, and DOI contain 2015; their relationship is unresolved here. The collection describes these reports as records of workshop discussions and open problems.
\item Lazarsfeld, \emph{Brill--Noether--Petri Without Degenerations}. Supplied literature reference: \url{https://doi.org/10.4310/JDG/1214440116}.
\item \emph{On Hyperplane Sections of K3 Surfaces}. Supplied literature reference: \url{https://arxiv.org/1507.05002}.
\item Haburcak, \emph{Curves on Brill--Noether Special K3 Surfaces}. Supplied literature reference: \url{https://arxiv.org/2306.11664}.
\item Literature assessment dated exactly \texttt{2026-08-21}: the supplied curation labels the problem open and states that the K3 locus has been studied using period maps, Wahl maps, and Brill--Noether theory, that Lazarsfeld's theorem has broad extensions, and that the two requested conclusions remain unsettled in the stated generality. This is supplied dated curation derived from the OWR record, not an independently verified literature assessment. Its time-sensitive status and cited claims require independent verification.
\item The registry labels the statement \texttt{corrected\_verified} and reports comparison with the official Oberwolfach report. This is a supplied verification claim, not verification performed for this draft.
\end{itemize}
\end{document}
